\begin{abstract}
  Coordination mechanisms such as barriers, atomic commit, and agreement protocols
  are ubiquitous in distributed systems, yet we lack a general semantic account of
  when coordination is \emph{intrinsically required} by a specification rather than
  introduced by a programming model or implementation strategy.
  We show that a simple monotonicity condition on which outcomes are
  observable (i.e., lie in $\Obs(H)$)
  exactly characterizes when a distributed specification admits a coordination-free,
  consistent implementation under asynchrony.

  Our criterion is phrased directly at the level of history-to-observation
  specifications and outcome orders, and is therefore independent of any particular
  programming model.
  Unlike CALM-style results, which formulate monotonicity at the level of input or
  program semantics, we place monotonicity on outcomes deemed observable by $\Obs$.
  This shift strictly generalizes program-level monotonicity: even specifications
  with non-monotone internal behavior may admit natural monotone observations and
  thus coordination-free implementations.

  We model executions as Lamport histories and define coordination-freedom as
  correctness under all causally consistent asynchronous executions, without
  excluding any admissible histories.
  The resulting \emph{Coordination Criterion} provides a \emph{sharp} semantic boundary
  (a dichotomy at the level of observable outcomes)
  between coordination-free and coordination-requiring specifications, and
  unifies classical results including CAP/linearizability, snapshots, agreement
  tasks, transactional isolation, and invariant confluence, with CALM arising as a
  special case.
\end{abstract}


\title{The Coordination Criterion}
\author{Joseph M. Hellerstein}
\affiliation{%
  \institution{UC Berkeley \& Amazon Web Services}
  \city{Berkeley, CA}
  \country{USA}
}
\date{}
\maketitle



% ================================================================
\section{Introduction}
\label{sec:introduction}
% ================================================================

Coordination is a recurring theme in distributed computing.
Protocols such as agreement, atomic commit, synchronous barriers, and termination
detection are widely used to ensure correct behavior under asynchrony and
partial failure.
At the same time, coordination is costly: it introduces latency, complexity,
and sensitivity to failures.
This motivates a fundamental question:

\begin{quote}
  When is coordination \emph{intrinsically necessary} to implement a distributed
  specification?
\end{quote}

% This question has appeared in several literatures.
% In the CAP theorem of Gilbert and Lynch~\cite{gilbert2002cap}, coordination manifests 
% as the unavailability required to maintain strong consistency during network partitions:
% to implement a linearizable read--write object, operations must block or abort
% when cross-partition agreement (e.g., quorum intersection) is impossible.
% The CALM principle~\cite{hellerstein2010calm} introduced the question above and 
% conjectured that monotonicity was the semantic differentiator;
% the CALM Theorem of Ameloot et al.~\cite{ameloot2013relational} 
% formalized the principle in the frame of logic languages for declarative networking and databases.
% In distributed computability theory, impossibility and hierarchy results
% characterize tasks that cannot be solved without waiting or stronger primitives
% \cite{herlihy1991waitfree,herlihy1999topological}.

% Each of these approaches addresses coordination under specific modeling
% assumptions and for particular classes of specifications.
% CAP-style results focus on a narrow family of strong consistency conditions
% such as linearizability and reason about availability under partitions, but do
% not provide a general semantic boundary between coordination-free and
% coordination-requiring specifications.
% CALM characterizes when relational transducer programs can be executed
% coordination-free, but is tied to a specific programming model and definition of 
% set-based monotonicity rather than arbitrary distributed specifications and outcome
% orders.
% Classical computability and hierarchy results analyze which tasks are solvable
% with particular primitives and progress conditions (e.g., wait-freedom), but
% do not separate the semantic necessity of coordination from the liveness and
% implementability assumptions of the underlying model.

The relationship between specification-level correctness requirements and
concrete coordination mechanisms has been explored in several lines of work.
In the CAP theorem~\cite{gilbert2002cap}, coordination appears as the loss of
availability required to maintain strong consistency under partitions, while
the CALM principle~\cite{hellerstein2010calm} connects coordination-freedom to
monotonicity in declarative programs.
Distributed computability theory similarly characterizes which tasks require
waiting or stronger primitives, most notably via object- and task-based
hierarchies such as consensus numbers~\cite{herlihy1991waitfree,herlihy1999topological}.
Each of these approaches isolates important aspects of coordination, but does so
under specific modeling assumptions and for particular classes of specifications.

What is missing is a general semantic account that explains coordination
\emph{as a property of the specification itself}, independent of programming
model, protocol structure, or particular consistency conditions.
Equivalently, we would like a single semantic boundary theorem that says
exactly when coordination is intrinsically required by a specification across
different models and implementation techniques.
Our Coordination Criterion provides such a sharp boundary theorem: it identifies a
monotonicity condition on specifications that is necessary and sufficient for
coordination-freedom. Because the criterion is phrased directly over observable
outcomes and histories, it applies uniformly across distributed objects,
agreement tasks, consistency models, and even optimization problems, with CALM-
and CAP-style results arising as special cases.

Historically, results on coordination in distributed systems have been developed
at different levels of abstraction, from low-level read/write objects and
consistency conditions to higher-level program and problem semantics.
Insights at one level have not always transferred cleanly to others.
By formulating coordination as a semantic property of specifications over
histories and observable outcomes, our framework is agnostic to
this choice of level, allowing coordination requirements to be analyzed
uniformly across classical object-level semantics and higher-level formulations.

Conceptually, the key to our results is to shift the focus from monotonicity
of input semantics (e.g., a query or update operator in a particular
programming model, as in CALM) to monotonicity of the specification map
with respect to an outcome order on observable outcomes.
Once we choose a possibility function $\Poss$, observability function $\Obs$, and
outcome order $\Ord$, our criterion asks only whether the induced specification
$\Spec = (\Poss, \Obs, \Ord)$ is future-monotone; any particular
programming model then
appears merely as an implementation vehicle for realizing such monotone
future-monotone specifications.

\paragraph{Contribution.}
We identify a general semantic boundary between
distributed specifications that admit coordination-free implementations and
those that do not.
Our contributions are as follows:
\begin{itemize}[nosep]
  \item \textbf{Semantic framework.} We formalize a minimalist, general semantic model
        for distributed specifications based on asynchronous histories, an
        observation function, and an outcome order capturing non-contradiction
        (Section~\ref{sec:obs-consistency}).
  \item \textbf{Coordination Criterion.} We prove that a specification admits a
      coordination-free consistent implementation if and only if it is future-monotone
      under the chosen outcome order (our main
        theorem).
  \item \textbf{Semantic grounding of classical results.} We instantiate the
        framework to semantically explain, and in some cases recover as direct
        corollaries, a range of
        classical results about CAP/linearizability, snapshots, agreement tasks,
        transactional isolation levels, and invariant confluence
        (Section~\ref{sec:applications} and Appendix~\ref{sec:additional-applications}).
\end{itemize}

Our focus in this paper is on safety-style semantic conditions on observable
behaviors—what we will broadly call \emph{consistency} in
Section~\ref{sec:obs-consistency}—not liveness.
As such, the Coordination Criterion characterizes when a specification admits a
coordination-free implementation whose observable behaviors satisfy a given
consistency requirement under all admissible asynchronous executions, including
executions with failures such as message loss, delay, and process crashes.
It is deliberately agnostic about progress guarantees such as termination or
wait-freedom, as well as the details of failure-recovery mechanisms, which are
orthogonal to the semantic questions studied here.
Accordingly, the Coordination Criterion focuses on when coordination is
semantically unavoidable, not on how non-monotonicity is resolved or how
progress is ensured.

% Our criterion is agnostic about how non-monotonicity is resolved—by
% protocol, interpretation, or additional semantic structure—focusing only on when
% such resolution is unavoidable.
% Unlike monotonicity results phrased within specific programming or replicated-object
% models, the Coordination Criterion gives a necessary and sufficient condition for
% coordination-freedom for specifications defined over standard Lamport histories.
% When we invoke classical impossibility or hierarchy theorems (e.g., for
% wait-free consensus and \(k\)-set agreement), we adopt their usual liveness
% assumptions, but these assumptions come from those external theorems rather
% than from our semantic notion of monotonicity.

Throughout, specifications in our framework are compositional and need not
correspond to entire distributed services.
They may describe the behavior of components, interfaces, or subroutines
within a larger system, allowing the Coordination Criterion to be applied
at multiple levels of specification.
The typical question is not whether an entire system uses coordination at all, but
which parts of its logic intrinsically require coordination and which
provably do not.

To make the discussion concrete, we will occasionally refer to a simple
running example: a single read--write register replicated at two processes
in an asynchronous message-passing system.
We will consider two specifications for this object over the same
computational model and observations: a highly available, eventually
consistent register and a linearizable register.
Throughout the paper we use this example to illustrate how choices of
observation, outcome order, and specification determine whether
coordination is semantically necessary.